\section{Trace byte cursors separately from buffered bits}
The worked stored stream occupies 46 bytes. Its first compressed byte is at offset
two, and its checksum begins at offset 42. The following offsets are relative to the
stream and output origins; helper argument cursors are separate from the retained
state shown here.

\Needspace{10\baselineskip}
\begin{longtable}{p{48mm}p{29mm}p{30mm}p{35mm}}
\toprule Boundary & Next input offset & Buffered bits & Produced output bytes \\\midrule
Initialized ranges & 2 & 0 & 0 \\
Three-bit header consumed & 3 & 5 & 0 \\
Stored alignment complete & 3 & 0 & 0 \\
LEN and NLEN consumed & 7 & 0 & 0 \\
Stored payload copied & 42 & 0 & 35 \\
Exact end checks passed & 42 & 0 & 35 \\\bottomrule
\end{longtable}
A cursor at offset three does not mean all preceding bits were consumed: the header
read leaves five bits in the reservoir. Stored alignment removes those buffered bits
without consuming another byte. Conversely, reaching the output end says nothing by
itself about the compressed cursor. In the extra-body fixture, output still reaches
35 while the compressed end is 43; the unmatched byte prevents acceptance.

The two-block fixture makes the final flag observable. Its first stored block produces
seven row bytes and has final flag zero. Continue through \texttt{block\_done}: the
next operation reads another block header rather than starting Adler verification.
Only the final block proceeds to exact end checks. Empty stored blocks can take this
same route without copying a byte.
